\documentclass[letterpaper, 10 pt, conference]{ieeeconf}
\IEEEoverridecommandlockouts
\overrideIEEEmargins

\usepackage{amsmath,amssymb,amsfonts}
\usepackage{graphicx}
\usepackage{cite}

\title{\bf Formal Guarantees for Commit-Time Governance: Viability, Barrier Certificates, and Reachability}
\author{[Authors anonymized for submission]}

\begin{document}
\maketitle
\thispagestyle{empty}
\pagestyle{empty}

%%%%%%%%%%%%%%%%%%%%%%%%%%%%%%%%%%%%%%%%%%%%%%%%%%%%%%%%%%%%%%%%%%%%%%%%%%%%%%%%
\begin{abstract}
We formalize the relationship between exact recoverability and the executable commit-time gate. Commit admissibility is expressed as a viability condition under lag, and the gate is shown to implement a conservative approximation using barrier certificates. We establish soundness: all states accepted by the gate belong to a certified subset of the robust recoverable region. Approximation error is characterized and shown to be localized near the boundary. The results demonstrate that the executable gate is a principled approximation of an exact reachability condition.
\end{abstract}

%%%%%%%%%%%%%%%%%%%%%%%%%%%%%%%%%%%%%%%%%%%%%%%%%%%%%%%%%%%%%%%%%%%%%%%%%%%%%%%%
\section{Introduction}

Prior work established commit-time admissibility and demonstrated an executable gate enforcing it. However, the gate operates on a conservative approximation of recoverability.

This paper formalizes the relationship between:
\begin{itemize}
\item exact recoverability (defined via lag-reachable sets),
\item viability (defined via infinite-horizon safety and recovery),
\item and the certified subset used in the implementation.
\end{itemize}

The goal is to show that the executable gate is not heuristic, but a sound approximation of a formally defined condition.

%%%%%%%%%%%%%%%%%%%%%%%%%%%%%%%%%%%%%%%%%%%%%%%%%%%%%%%%%%%%%%%%%%%%%%%%%%%%%%%%
\section{Viability Formulation}

We consider dynamics:
\begin{equation}
\dot{\mathbf{x}} = f(\mathbf{x}, u, d)
\end{equation}

with bounded disturbance $d \in \mathcal{D}$.

\begin{definition}[Viability Kernel]
\begin{equation}
\text{Viab}(\tau) =
\left\{
\mathbf{x}_0 :
\forall d(\cdot), \exists u(\cdot),\;
\mathbf{x}(t) \notin \mathcal{C} \;\forall t \in [0,\tau],\;
\mathbf{x}(\tau) \in \mathcal{R}(0)
\right\}
\end{equation}
\end{definition}

This captures states that remain safe during lag and are recoverable afterward.

\begin{theorem}[Viability–Recoverability Equivalence]
\begin{equation}
\text{Viab}(\tau) = \mathcal{R}_{rob}(\tau)
\end{equation}
\end{theorem}

\textit{Argument.}
Both sets require that all disturbance realizations over $[0,\tau]$ avoid collapse and lead to states from which recovery is possible. The definitions coincide.

%%%%%%%%%%%%%%%%%%%%%%%%%%%%%%%%%%%%%%%%%%%%%%%%%%%%%%%%%%%%%%%%%%%%%%%%%%%%%%%%
\section{Barrier Certificates}

We define barrier functions:
\begin{align}
B_F(\tilde{\mathbf{x}}) &= d_F(\tilde{\mathbf{x}}) - \epsilon_{crit}(\tau) \\
B_D(\tilde{\mathbf{x}}) &= d_{max}(\tau) - d_\Delta(\tilde{\mathbf{x}})
\end{align}

\begin{assumption}[C1: Face Safety]
If $B_F > 0$, then all lag-reachable trajectories remain separated from $\mathcal{C}$.
\end{assumption}

\begin{assumption}[C2: Recoverability]
If $B_D > 0$, then there exists a control law driving the system to $\mathcal{A}$.
\end{assumption}

\begin{definition}[Certified Safe Subset]
\begin{equation}
\mathcal{S}_{cert}(\tau) =
\left\{
\mathbf{x} :
B_F(\tilde{\mathbf{x}}) > 0
\;\wedge\;
B_D(\tilde{\mathbf{x}}) > 0
\right\}
\end{equation}
\end{definition}

%%%%%%%%%%%%%%%%%%%%%%%%%%%%%%%%%%%%%%%%%%%%%%%%%%%%%%%%%%%%%%%%%%%%%%%%%%%%%%%%
\section{Soundness}

\begin{theorem}[Certified Subset Soundness]
Under assumptions C1--C2:
\begin{equation}
\mathcal{S}_{cert}(\tau) \subseteq \mathcal{R}_{rob}(\tau)
\end{equation}
\end{theorem}

\textit{Proof.}
If $B_F > 0$, then trajectories remain separated from collapse during lag.  
If $B_D > 0$, then recovery to $\mathcal{A}$ exists.  
Thus both conditions imply membership in $\mathcal{R}_{rob}(\tau)$. \hfill $\blacksquare$

\begin{corollary}[Gate Soundness]
\begin{equation}
\text{ALLOW}_{Paper6}(u;\mathbf{x},\tau)
\;\Rightarrow\;
\text{ALLOW}_{exact}(u;\mathbf{x},\tau)
\end{equation}
\end{corollary}

%%%%%%%%%%%%%%%%%%%%%%%%%%%%%%%%%%%%%%%%%%%%%%%%%%%%%%%%%%%%%%%%%%%%%%%%%%%%%%%%
\section{Approximation Error}

\begin{definition}[Error Set]
\begin{equation}
\mathcal{E}(\tau) =
\mathcal{R}_{rob}(\tau) \setminus \mathcal{S}_{cert}(\tau)
\end{equation}
\end{definition}

\begin{proposition}[Error Localization]
The error set is concentrated near the boundary of $\mathcal{R}_{rob}(\tau)$.
\end{proposition}

\textit{Interpretation.}
States close to collapse or far from equilibrium fail the conservative test despite being recoverable.

\begin{proposition}[No False Positives]
The gate admits no states outside $\mathcal{R}_{rob}(\tau)$.
\end{proposition}

%%%%%%%%%%%%%%%%%%%%%%%%%%%%%%%%%%%%%%%%%%%%%%%%%%%%%%%%%%%%%%%%%%%%%%%%%%%%%%%%
\section{Discussion}

The executable gate is a conservative approximation of a viability condition. The approximation trades completeness for safety, ensuring that only states with certified recoverability are accepted.

%%%%%%%%%%%%%%%%%%%%%%%%%%%%%%%%%%%%%%%%%%%%%%%%%%%%%%%%%%%%%%%%%%%%%%%%%%%%%%%%
\section{Conclusion}

We established that the commit-time gate implements a sound approximation of recoverability. The connection to viability theory provides a formal basis for the implementation and clarifies the nature of approximation error.

\end{document}
